\item Un \textit{calorímetro de flujo} es un aparato que se usa para medir el calor específico de un líquido. La técnica de calorimetría de flujo involucra la medición de la diferencia de temperatura entre los puntos de entrada y salida de una corriente de líquido que circula mientras se agrega energía por calor a una rapidez conocida. Un líquido de densidad $\rho$ fluye a través del calorímetro con una rapidez de flujo volumétrico $R$. En estado estacionario, se establece una diferencia de temperatura $\Delta T$ entre los puntos de entrada y salida cuando se proporciona energía a una rapidez $P$. ¿Cuál es el calor específico del líquido?
